\emph{Setting and sign.} The derived metric is $dT^{2} - \sum_{a}T^{2P_{a}}(dx^{a})^{2}$ with $\sum_{a}P_{a} = 1$, so $\sqrt{|g|} = T$. The linearised computation is done in the overall opposite sign $-dT^{2} + \sum_{a}T^{2P_{a}}(dx^{a})^{2}$, which leaves the Ricci tensor and every mode equation unchanged. The points used are on $K_{5}$ ($\sum P = \sum P^{2} = 1$) and on the disk; the reduction of Theorem \claimref{synthesis.scalar-kasner-wall-is-reduced-five-dimensional-vacuum-wall}, $p_{i} = (P_{i} + P_{4}/2)/k$, $k = 1 + P_{4}/2$, $c^{2} = 3P_{4}^{2}/(2k^{2})$, sends both $(\tfrac12, \tfrac12, \tfrac12, -\tfrac12)$ and $(0, 0, 0, 1)$ to $(\tfrac13, \tfrac13, \tfrac13)$ with $c^{2} = \tfrac23$, and a generic point such as $(\tfrac15, \tfrac25, \tfrac45, -\tfrac25)$ to an anisotropic disk point. Equal observed exponents $p$ on $K_{5}$ require $3p + P_{4} = 1$ and $3p^{2} + P_{4}^{2} = 1$, so $12p^{2} - 6p = 0$: $p = 0$ or $p = \tfrac12$, the two isotropic points. The vacuum circle is $p(u) = (-u, 1 + u, u(1 + u))/(1 + u + u^{2})$ (Theorem \claimref{synthesis.bkl-bounce-invariant-holonomy-and-bound}); the Schwarzschild interior $(-\tfrac13, \tfrac23, \tfrac23)$ is $u = 1$ (Lemma \claimref{interior.kasner-form-interior}); the hidden direction is that of Proposition \claimref{synthesis.unit-ball-clauses-hold-at-one-epoch-unless-hidden-direction-co-expands}, read at its zero mode.

\emph{(i)} Perturb the metric by $h_{ab} = h_{ba} = T^{2P_{b}}A(T)\cos(kx^{j})$, $a, b, j$ distinct. To first order the inverse metric gains $g^{ab} = -h_{ab}T^{-2P_{a} - 2P_{b}}$. The Christoffel symbols that change are those with one index $T$ or $j$ and the pair $\{a, b\}$; collecting the terms of $\delta R_{ab}$,
\[
\delta R_{ab} = \tfrac12 T^{2P_{b}}\Bigl[\ddot A + \Bigl(\sum_{c}P_{c} - 2P_{a} + 2P_{b}\Bigr)\frac{\dot A}{T} + 2P_{b}\Bigl(\sum_{c}P_{c} - 1\Bigr)\frac{A}{T^{2}} + k^{2}T^{-2P_{j}}A\Bigr]\cos(kx^{j}),
\]
which with $\sum_{c}P_{c} = 1$ is the stated operator times $\tfrac12 T^{2P_{b}}\cos(kx^{j})$; all other components of $\delta R$ vanish to first order (the perturbation is transverse, $k$ along $j$ and the polarisation in the $ab$ plane, and traceless). The general-$P$ form was computed symbolically; the check recomputes the full Ricci tensor to first order exactly, with a formal engine on terms $cT^{n}A^{(d)}\cos$ or $\sin(kx^{j})\,k^{m}$, at the two isotropic points and the generic point of $K_{5}$ (three choices of $a, b, j$ there), two points of the vacuum circle including the Schwarzschild interior, two four-dimensional disk points and two five-dimensional stiff points of the segment of (v), and finds the stated equation and no other first-order component; a control with $P_{a}$ and $P_{b}$ exchanged in the friction fails. At the vacuum points the background is Ricci flat, so $\delta R_{ab} = 0$ is the field equation. At a disk point the field equation is $\Ric = -8\pi G\,d\phi\otimes d\phi$ (Theorem \claimref{synthesis.scalar-kasner-wall-is-reduced-five-dimensional-vacuum-wall}; stiff matter at rest is the same source with $\rho_{0} = \tfrac12\dot\phi^{2}$, Theorem \claimref{synthesis.kasner-wall-matter-maxwell-bounces-stiff-moves-residue}(iii)) with $\phi$ a function of $T$ alone, so its $ab$ component has no first-order source, and the scalar's own equation is not sourced either: $h^{ab}\partial_{a}\partial_{b}\phi = 0$ and the background $\Gamma^{T}_{ab}$ vanishes for $a \ne b$. The check confirms that at the disk points the background Ricci tensor has only its $TT$ component. In the variable $\eta_{j} = T^{1 - P_{j}}/(1 - P_{j})$ the equation is $(z^{2}A')' + k^{2}z^{2}A = 0$ with $z^{2} = T^{1 + 2P_{b} - 2P_{a} - P_{j}}$, a pump exponent $w = 2(P_{b} - P_{a})$ in the notation of (ii), so $\nu = |w|/(2(1 - P_{j})) = |P_{a} - P_{b}|/(1 - P_{j})$. For $b$ along the hidden direction, $g_{a4} = b^{2}A_{a}$ with $b = T^{P_{4}}$ is the Kaluza-Klein ansatz (Fact \claimref{synthesis.imported-kaluza-klein-ansatz}), so $A$ is the Kaluza-Klein vector's component; for $a, b$ observed it is the cross polarisation of a gravitational wave travelling along $j$.

\emph{(ii)} Put $q_{j} = 1 - p_{j}$, $s = k\tau^{q_{j}}/q_{j}$, $\alpha = w/(2q_{j})$ and $\chi = s^{-\alpha}Z(s)$. With $d/d\tau = k\tau^{-p_{j}}d/ds$ and $\tau^{-2} = k^{2}\tau^{-2p_{j}}/(q_{j}^{2}s^{2})$, the operator $\ddot\chi + (1 + w)\dot\chi/\tau + (k^{2}\tau^{-2p_{j}} - \xi c^{2}\tau^{-2})\chi$ equals $k^{2}\tau^{-2p_{j}}s^{-\alpha}\bigl(Z'' + Z'/s + (1 - \nu^{2}/s^{2})Z\bigr)$ with $\nu^{2} = \alpha^{2} + \xi c^{2}/q_{j}^{2}$: on $Z = s^{n}$ the coefficient of $s^{n - 2}$ is $e^{2} + we - \xi c^{2}$ with $e = q_{j}(n - \alpha)$, and $e^{2} + we = q_{j}^{2}(n^{2} - \alpha^{2})$. The check verifies the identity on $Z = s^{n}$ for six values of $n$ and $120$ choices of $(w, \xi, c^{2}, p_{j})$, with a control at a shifted order. On the disk the Ricci scalar is $R\tau^{2} = \sum p^{2} - 1 = -c^{2}$ in the mostly-minus-time sign of the computation (the check computes it from the engine at three disk points); the sign of the $\xi$ term is the one for which $\xi = \tfrac16$ is the conformal coupling: at $(\tfrac13, \tfrac13, \tfrac13)$ the metric is conformally flat with $a \propto \eta^{1/2}$, a conformally coupled massless field has $u = a\chi$ obeying $u'' + k^{2}u = 0$, order $\tfrac12$, and the formula gives $\nu^{2} = \tfrac16\cdot\tfrac23\cdot\tfrac94 = \tfrac14$, as checked.

\emph{(iii)} The momentum conjugate to $\chi$ is $\tau^{1 + w}\dot\chi$ up to a $k$-independent factor, so the mode in its adiabatic vacuum before exit is $\chi = Ns^{-\alpha}H^{(1)}_{\nu}(s)$ with the canonical condition $\tau^{1 + w}|\chi\dot\chi^{*} - \chi^{*}\dot\chi| = 1$. With $ds/d\tau = q_{j}s/\tau$ and $W[H^{(1)}_{\nu}, H^{(2)}_{\nu}](s) = -4i/(\pi s)$ this reads $|N|^{2}q_{j}\,\tau^{w}s^{-2\alpha}\cdot 4/\pi = 1$; since $2\alpha q_{j} = w$, $\tau^{w}s^{-2\alpha} = (q_{j}/k)^{2\alpha}$, so $|N|^{2} = \tfrac{\pi}{4q_{j}}(k/q_{j})^{2\alpha}$ and
\[
|\chi_{k}|^{2} = |N|^{2}s^{-2\alpha}|H^{(1)}_{\nu}(s)|^{2} = \frac{\pi}{4q_{j}}\,\tau^{-w}\,|H^{(1)}_{\nu}(s)|^{2},
\]
in which $k$ enters only through $s$, as for order $0$ in Theorem \claimref{synthesis.kasner-approach-gives-horizon-exit-and-index-three}(iii). For $s \to 0$, $|H^{(1)}_{\nu}(s)| \sim \Gamma(\nu)(s/2)^{-\nu}/\pi$ for $\nu > 0$ (Fact \claimref{synthesis.imported-bessel-reduction-and-small-argument-hankel-forms}), so at fixed $\tau$, $\Delta^{2} = k^{3}|\chi_{k}|^{2}/(2\pi^{2}) \propto k^{3 - 2\nu}$ and $n - 1 = 3 - 2\nu$. For the off-diagonal modes $\chi = A$; the transverse component is unchanged by linearised diffeomorphisms (for the Kaluza-Klein vector, $A_{a} \mapsto A_{a} + \partial_{a}\lambda$ with $\lambda \propto \cos(kx^{j})$ leaves $A_{a}$, $a \ne j$, fixed), so the index is that of a gauge-invariant quantity, but for a vector the field strength is the physical field: its magnetic part $F_{ja} = \partial_{j}A_{a}$ carries one more power of $k$, index $5 - 2\nu$. Once the hidden scale stops moving the Kaluza-Klein vector is a four-dimensional Maxwell field, conformal, with energy density per logarithmic interval of $k$ proportional to $k^{2}\Delta_{A}^{2}$ over the fourth power of the scale factor, so its energy spectrum has the field-strength index $5 - 2\nu$, not the potential's (unlike a graviton, whose strain is the physical amplitude). The electric part comes from the time derivative; where the leading superhorizon branch $s^{-\alpha - \nu}$ is constant in time ($\alpha = -\nu$, the Kaluza-Klein vector for $p > \tfrac14$) it is set by the next branch and is bluer still (argued). A vector sources an isotropic curvature perturbation only quadratically, through $F^{2}$, and a local conversion of a Gaussian field does not lower its index (\claimref{synthesis.local-conversion-cannot-redden-a-gaussian-spectator}).

\emph{(iv)} On the vacuum circle, with $d = 1 + u + u^{2}$: $1 - p_{1} = (1 + u)^{2}/d$ and $p_{2} - p_{3} = (1 - u^{2})/d$, so $\nu_{1} = |1 - u|/(1 + u)$; $1 - p_{2} = u^{2}/d$ and $|p_{1} - p_{3}| = u(2 + u)/d$, so $\nu_{2} = (2 + u)/u$; $1 - p_{3} = 1/d$ and $|p_{1} - p_{2}| = (1 + 2u)/d$, so $\nu_{3} = 1 + 2u$. The indices are $3 - 2|1 - u|/(1 + u) \in (1, 3]$, $1 - 4/u \in [-3, 1)$ and $1 - 4u \le -3$; since $|1 - u|/(1 + u) < 1$ the first exceeds $1$, so the largest minus the smallest exceeds $1 - (1 - 4u) = 4u \ge 4$; the check verifies the three orders and the spread exactly at $u = 1, \tfrac32, 2, \tfrac73, 4, 10$. For a red index along one axis the spectrum along another axis differs by a power at least $k^{4}$; a statistically isotropic spectrum, as measured (Fact \claimref{synthesis.imported-planck-2018-isocurvature-and-isotropy-bounds}), has the same index along every direction. At a disk point a curvature-coupled scalar has $\nu_{j} = \sqrt{\xi}|c|/(1 - p_{j})$, again equal on the observed axes only at equal exponents when nonzero.

\emph{(v), vacuum ends.} At the isotropic points $P_{a} = P_{b}$ for observed $a, b$, so the cross modes have $\nu = 0$ and the stored index; for the Kaluza-Klein vector $\nu = |P_{a} - P_{4}|/(1 - P_{j})$ is $1/\tfrac12 = 2$ at $(\tfrac12, \tfrac12, \tfrac12, -\tfrac12)$ and $1/1 = 1$ at $(0, 0, 0, 1)$, for all six choices of $a \ne j$ (checked), giving $3 - 4 = -1$ and $3 - 2 = 1$. A conformally coupled scalar or a four-dimensional gauge vector without a modulus coupling has order $\tfrac12$ there (conformal invariance at a conformally flat point), index $2$.

\emph{(v), the segment.} For $P = (p, p, p, 1 - 3p)$, $\sum P^{2} = 1 + 12p^{2} - 6p$, which is at most $1$ exactly for $0 \le p \le \tfrac12$ and equals $1$ only at the ends; the interior points are those of the five-dimensional analogue of Theorem \claimref{synthesis.kasner-wall-matter-maxwell-bounces-stiff-moves-residue}(iii): stiff matter in five dimensions is a constant potential that puts the weights on $\sum P = 1$, $\sum P^{2} < 1$ (standard, cited from memory; not in the store, and no stored field supplies five-dimensional stiff matter), so only the vacuum ends are reached by the store's own fields. The cross modes of observed pairs have $\nu = 0$ there; the ratio $|P_{h}|/(1 - p) = |1 - 3p|/(1 - p)$ is at most $1$ on the segment ($1$ minus it is $2p/(1 - p)$ for $p \le \tfrac13$ and $(2 - 4p)/(1 - p)$ for $p \ge \tfrac13$). The Kaluza-Klein vector has $\nu = |P_{h} - p|/(1 - p) = |1 - 4p|/(1 - p)$, in $[0, 2]$, zero at $p = \tfrac14$; for $p \ge \tfrac14$ its index is $3 - 2(4p - 1)/(1 - p) = (5 - 11p)/(1 - p)$, decreasing from $3$ to $-1$ at $p = \tfrac12$, and for $p \le \tfrac14$ it is $(1 + 5p)/(1 - p) \ge 1$. The check verifies these on a grid of $61$ points of the segment (and the disk membership on $181$ points of $[-\tfrac12, 1]$).

\emph{(vi)} The list compared was fixed in advance: the cross modes, a conformal scalar or gauge vector, the Kaluza-Klein vector at both isotropic points, and a curvature-coupled scalar at the minimal ($\xi = 0$, the minimal coupling of Fact \claimref{synthesis.imported-einstein-cartan-matter-coupling}) and conformal ($\xi = \tfrac16$) values: about eight values. With $n_{s} - 1 = -0.0351 \pm 0.0042$ (Fact \claimref{synthesis.imported-planck-2018-lcdm-parameters}) the three-standard-deviation window is $[-0.0477, -0.0225]$; the values $2$ or more, $1$ and $-1$ are outside it, and $|-1 + 0.0351| = 0.9649 > 229 \times 0.0042 = 0.9618$. For the scalar at the isotropic point, $n - 1 = 3 - 2\sqrt{3\xi/2} = 3 - \sqrt{6\xi}$, so $\xi = (3 - (n - 1))^{2}/6$, which over the window runs from $3.0225^{2}/6 \in (1.5225, 1.5226)$ to $3.0477^{2}/6 \in (1.5480, 1.5481)$; $\xi = \tfrac32$ gives $\nu = \tfrac32$ and $n = 1$ exactly. No entity of the store fixes a curvature coupling in this range; the free coupling is one parameter fitted to one number. The segment itself was added to the comparison after the vacuum ends had been compared; along it the Kaluza-Klein vector's index $x = (5 - 11p)/(1 - p)$ inverts to $p = (5 - x)/(11 - x)$, decreasing in $x$, so the window $x \in [-0.0477, -0.0225]$ is met exactly for $p \in [5.0225/11.0225, 5.0477/11.0477] \subset (0.4556, 0.4570)$, interior to $(\tfrac13, \tfrac12)$. There $\nu = (3 - x)/2$, so the field strength's index $5 - 2\nu = 2 + x$ lies in $[1.9523, 1.9775]$: blue. On the whole segment $5 - 2\nu \ge 5 - 4 = 1$. So the potential's passage through the window is a property of the potential and not a route to the measured tilt; only the curvature coupling, one free number fitted to one target, matches it. In $\Cl(1,4)$ ($e_{0}^{2} = 1$, $e_{a}^{2} = -1$), $B_{5} = \sum_{a}P_{a}e_{0}e_{a}$ has $B_{5}^{2} = \sum P_{a}^{2} = 1$ on $K_{5}$, and $e_{a}B_{5} = P_{a}e_{a}e_{0}e_{a} = P_{a}e_{0}$, $B_{5}e_{a} = P_{a}e_{0}e_{a}e_{a} = -P_{a}e_{0}$ for the $a$ term (the other terms commute with $e_{a}$ up to sign and cancel in the difference), so $\tfrac12(e_{a}B_{5} - B_{5}e_{a}) = P_{a}e_{0}$, checked at three points: the order's numerator is the difference of these weights.
